\documentclass[11pt]{article}
\usepackage[margin=1in]{geometry}
\usepackage{amsmath,amssymb,amsthm}
\usepackage{booktabs}
\usepackage{bm}
\usepackage{tikz}
\usetikzlibrary{quantikz2}
\usepackage[colorlinks=true,linkcolor=blue!50!black,citecolor=blue!50!black,urlcolor=blue!50!black]{hyperref}

\newtheorem{theorem}{Theorem}
\newtheorem{lemma}[theorem]{Lemma}
\newtheorem{corollary}[theorem]{Corollary}
\theoremstyle{remark}
\newtheorem{remark}[theorem]{Remark}

\newcommand{\CX}{\mathrm{CX}}
\newcommand{\CZ}{\mathrm{CZ}}
\newcommand{\Pf}{\operatorname{Pf}}
\newcommand{\Arg}{\operatorname{Arg}}

\title{One more fold per node: block-ZXZ synthesis with\\ $\tfrac{21}{48}4^n$ CNOT gates}
\author{}
\date{}

\begin{document}
\maketitle

\begin{abstract}
The block-ZXZ decomposition of Krol and Al-Ars synthesizes an arbitrary $n$-qubit unitary with
$\tfrac{22}{48}4^n-\tfrac32 2^n+\tfrac53$ CNOT gates, the best known count for an exact analytic
method. We show that every recursion node admits one further CNOT cancellation. The free
single-qubit gates on the top qubit at each node can always be chosen so that the input-side
multiplexed rotation folds two CNOTs into the middle factor instead of one. A parity rule for
one angle and the intermediate value theorem for the other prove the choice exists. The result
is $\tfrac{21}{48}4^n-\tfrac32 2^n+2$ CNOTs, i.e.\ $18$ for a general three-qubit unitary
(previously $19$) and $(4^{n-2}-1)/3$ fewer than block-ZXZ for every $n\ge 3$.
\end{abstract}

\section{Context}

For exact synthesis of a generic $U\in U(2^n)$ from CNOTs and single-qubit gates, parameter
counting gives the lower bound $\lceil(4^n-3n-1)/4\rceil\approx\tfrac{12}{48}4^n$~\cite{smb}. The
best analytic constructions are recursive. The quantum Shannon decomposition (QSD) with its
optimizations uses $\tfrac{23}{48}4^n-\tfrac32 2^n+\tfrac43$ CNOTs~\cite{smb}. The block-ZXZ
decomposition~\cite{zxz} lowers the leading coefficient to $\tfrac{22}{48}$. Numerical
optimization does better on small registers: for instance $14$ CNOTs for three qubits, the lower
bound, with brick-wall templates~\cite{brickwall}, and near-optimal counts via gradient
methods~\cite{squander}. But no formula or guarantee is known for these, and universality
of the lower-bound templates is conjectural. For analytic methods, the leading coefficient had
moved from $\tfrac{23}{48}$ (2004) to $\tfrac{22}{48}$ (2024). This note moves it to $\tfrac{21}{48}$.

\section{Block-ZXZ recap}

Let $n\ge3$, $D=2^{n-1}$, let $t=q_{n-1}$ be the top qubit (most significant) and write
$U=\begin{psmallmatrix}X&Y\\ \ast&\ast\end{psmallmatrix}$ in $D\times D$ blocks. Block-ZXZ
factors~\cite{zxz}
\begin{equation}\label{eq:zxz}
U=(A_1\oplus A_2)\,(H\otimes I)\,(I\oplus B)\,(H\otimes I)\,(I\oplus C),\qquad
C=-i\,U_X^\dagger U_Y ,
\end{equation}
where $X=S_XU_X$, $Y=S_YU_Y$ are polar decompositions. The factors $A_1\oplus A_2$ and $I\oplus C$
are demultiplexed, e.g.\ $I\oplus C=(I\otimes V_c)\,\Delta(\theta)\,(I\otimes W_c)$. Here
$\Delta(\theta)=\sum_j|j\rangle\langle j|\otimes R_z(\theta_j)$ is an $R_z$ rotation on $t$
multiplexed by the lower register, and $e^{-i\theta_j}$ runs over the eigenvalues of $C^\dagger$
(so $\theta_j\equiv\mu_j \bmod 2\pi$ for the eigenphases $\mu_j$ of $C$). The same is done for the
resulting middle factor. A multiplexor with $n-1$ controls costs $D$ CNOTs in Gray-code
form~\cite{mottonen}. Block-ZXZ saves one CNOT on each outer multiplexor: its last CNOT
$\CX(q,t)$ meets the Hadamard, $(H\otimes I)\CX(q,t)=\CZ(q,t)(H\otimes I)$, and
$\CZ(q,t)=I\oplus Z_q$ is absorbed into the block-diagonal middle factor. Recursing, with
two-qubit leaves of $3$ CNOTs (exact) or $2$ (up to a diagonal that migrates into the next leaf),
gives
\[
c_{\rm ZXZ}(n)=4c_{\rm ZXZ}(n-1)+3\cdot 2^{n-1}-5 .
\]

\begin{figure}[t]
\centering
\begin{quantikz}[column sep=0.22cm, row sep={0.75cm,between origins}]
\lstick{$q_2$} & \gate{u^\dagger} & & \gate[3]{\Delta_c} & & \gate{H} & \gate[3]{\Delta_b} & \gate{H} & & \gate[3]{\Delta_a} & & \\
\lstick{$q_1$} & & \gate[2]{W_c} & & \gate[2]{W_b} & & & & \gate[2]{V_b} & & \gate[2]{V_a} & \\
\lstick{$q_0$} & & & & & & & & & & &
\end{quantikz}
\caption{A block-ZXZ node for $n=3$ (time runs left to right). $W,V$ are two-qubit leaves;
$\Delta_c,\Delta_b,\Delta_a$ are $R_z$ rotations on $q_2$ multiplexed by $q_1q_0$. Block-ZXZ
spends $2{+}3{+}2{+}4{+}2{+}3{+}3=19$ CNOTs. This work adds the free single-qubit gate $u^\dagger$
on the top qubit and lowers $\Delta_c$ from $3$ to $2$ CNOTs.}
\label{fig:node}
\end{figure}

\section{The extra fold}

In Gray-code order, $\Delta(\theta)$ ends with $\CX(q_0,t)\,R_z(\varphi)\,\CX(q_{n-2},t)$, where
\begin{equation}\label{eq:phi}
\varphi=\frac1D\sum_{j}(-1)^{j_{n-2}}\,\theta_j
\end{equation}
is the Walsh coefficient of the highest control bit. If $\varphi=0$, the two CNOTs are adjacent
and \emph{both} fold (Fig.~\ref{fig:fold}):
$(H\otimes I)\CX(q_0,t)\CX(q_{n-2},t)=\CZ(q_0,t)\CZ(q_{n-2},t)(H\otimes I)$ and
$\CZ(q_0,t)\CZ(q_{n-2},t)=I\oplus Z_{q_0}Z_{q_{n-2}}$ is again absorbed into the middle factor,
which becomes $(\cdots)(I\oplus B)(V_c\oplus V_cZ_{q_0}Z_{q_{n-2}})$ before it is demultiplexed.

\begin{figure}[t]
\centering
\textbf{(a) block-ZXZ}\\[2pt]
\begin{quantikz}[column sep=0.18cm, row sep={0.65cm,between origins}]
\lstick{$q_2$} & \gate{R_z(\varphi_0)} & \targ{} & \gate{R_z(\varphi_1)} & \targ{} & \gate{R_z(\varphi_2)} & \targ{} & \gate{R_z(\varphi)} & \targ{}\gategroup[3,steps=2,style={dashed,rounded corners,fill=blue!10,inner xsep=2pt},background,label style={label position=above,anchor=south,yshift=-0.1cm}]{{\footnotesize folds to $\CZ$}} & \gate{H} & \\
\lstick{$q_1$} & & & & \ctrl{-1} & & & & \ctrl{-1} & & \\
\lstick{$q_0$} & & \ctrl{-2} & & & & \ctrl{-2} & & & &
\end{quantikz}\\[8pt]
\textbf{(b) this work: choose $u$ so that $\varphi=0$}\\[2pt]
\begin{quantikz}[column sep=0.18cm, row sep={0.65cm,between origins}]
\lstick{$q_2$} & \gate{R_z(\varphi_0)} & \targ{} & \gate{R_z(\varphi_1)} & \targ{} & \gate{R_z(\varphi_2)} & \targ{}\gategroup[3,steps=3,style={dashed,rounded corners,fill=blue!10,inner xsep=2pt},background,label style={label position=above,anchor=south,yshift=-0.1cm}]{{\footnotesize both fold}} & \targ{} & \gate{H} & \midstick[3,brackets=none]{=} & \gate{H} & \ctrl{2} & \ctrl{1} & \\
\lstick{$q_1$} & & & & \ctrl{-1} & & & \ctrl{-1} & & & & & \control{} & \\
\lstick{$q_0$} & & \ctrl{-2} & & & & \ctrl{-2} & & & & & \control{} & &
\end{quantikz}
\caption{The input-side multiplexor $\Delta_c$ followed by the Hadamard of the middle factor
($n=3$). (a)~Block-ZXZ folds the last CNOT. (b)~If the Walsh coefficient $\varphi$ vanishes,
the last two CNOTs are adjacent, both become CZs, and $\CZ(q_0,q_2)\CZ(q_1,q_2)=I\oplus Z_0Z_1$
is absorbed into the next block-diagonal factor: $\Delta_c$ costs $D-2$ CNOTs instead of $D-1$.}
\label{fig:fold}
\end{figure}

\begin{lemma}[fold condition]\label{lem:fold}
The eigenvalue order and square-root branches in the demultiplexing of $I\oplus C$ can be chosen
so that $\varphi=0$ if and only if the eigenphases $\mu_1,\dots,\mu_D$ of $C$ admit a balanced
split, $|S|=D/2$, with
\begin{equation}\label{eq:split}
\textstyle\sum_{k\in S}\mu_k-\sum_{k\notin S}\mu_k\in 2\pi\mathbb Z .
\end{equation}
\end{lemma}
\begin{proof}
The order assigns eigenvalues to lower basis states $j$. Put $S$ on the states with
$j_{n-2}=0$; then~\eqref{eq:phi} reads $\varphi=\frac1D(\sum_S\theta_k-\sum_{S^c}\theta_k)$.
Since $\theta_k\in\mu_k+2\pi\mathbb Z$ can be shifted freely (a $2\pi$ shift is a sign on one
eigenvector, absorbed by $V_c,W_c$), $\varphi$ can be made $0$ iff~\eqref{eq:split} holds.
\end{proof}

A generic $C$ has no such split. The node offers free parameters, though: a single-qubit gate
$u$ on $t$ before the node (implemented as $u^\dagger$ in the circuit, Fig.~\ref{fig:node}) replaces
$U$ by $Uu$ and changes $C$. It costs no CNOTs.

\section{Existence of the gate $u$}

\paragraph{A sign-changing test function.}
Let $V\in U(D)$, $D=2^{n-1}\ge4$, and $r$ a $D$-th root of $\det V$. The exterior power
$R=\Lambda^{D/2}(V/r)$ acts on $\Lambda^{D/2}\mathbb C^D$, of dimension $N=\binom{D}{D/2}$. Since
$D/2$ is even, the wedge pairing on $\Lambda^{D/2}$ is symmetric. Together with the Hermitian
form, it defines a real structure preserved by $SU(D)$, so in a fixed real orthonormal basis
$R\in SO(N)$. Put
\[
P(V,r)=\Pf(R-R^{\mathsf T})=\prod_{\ell}2\sin\psi_\ell ,
\]
where $\pm\psi_\ell$ are the rotation angles of $R$. The eigenvalues of $R$ are the products
$\lambda_S=\prod_{k\in S}\lambda_k/r^{D/2}$ over half-subsets $S$, and $\lambda_S\lambda_{S^c}=1$.
Hence $P=0$ iff some $\lambda_S=\pm1$, iff $\lambda_S=\lambda_{S^c}$, iff~\eqref{eq:split} holds.
Two facts are used below. $N/2=\binom{D-1}{D/2-1}$ is odd (Lucas' theorem, since $D-1$ is all ones
in binary), so
$\Pf(R^{\mathsf T}-R)=\Pf(-(R-R^{\mathsf T}))=-P$. And replacing $r$ by $re^{2\pi i m/D}$
multiplies $R$ by $e^{-\pi i m}=(-1)^m$.

\paragraph{An involution.}
Let $u(\beta)=R_z(\alpha)R_y(\beta)$ on $t$. With $c=\cos\frac\beta2$, $s=\sin\frac\beta2$, the
top-row blocks of $Uu(\beta)$ are, up to the phases $e^{\mp i\alpha/2}$,
$X(\beta)=cX+sY$ and $Y(\beta)=-sX+cY$. At $\beta=\pi$ they are $(Y,-X)$, so
$U_{X(\pi)}=U_Y$, $U_{Y(\pi)}=-U_X$ and
\begin{equation}\label{eq:inv}
C(\pi)=i\,U_Y^\dagger U_X=C(0)^{-1}.
\end{equation}

\begin{lemma}[parity]\label{lem:parity}
Let $\tau_1,\dots,\tau_D$ be the eigenvalues of $X^{-1}Y$, none real (after the phases above,
$\tau_k\mapsto e^{i\alpha}\tau_k$). Along $\beta\in[0,\pi]$ the polar factors, hence $C(\beta)$,
are continuous. Take the continuous root $r(\beta)$ of $\det C(\beta)$ with
$r(0)=e^{i\sum_k\Arg\tau_k/D}$ and set $m=\#\{k:\operatorname{Im}\tau_k>0\}-D/2$. Then
$r(\pi)=\overline{r(0)}\,e^{2\pi i m/D}$ and
\[
P\big(C(\pi),r(\pi)\big)=(-1)^{m+1}\,P\big(C(0),r(0)\big).
\]
\end{lemma}
\begin{proof}
$\det X(\beta)=\det X\prod_k(c+s\tau_k)$ and $\det Y(\beta)=\det Y\prod_k(c-s/\tau_k)$ never vanish
because no $\tau_k$ is real. As $(c,s)$ moves from $(1,0)$ to $(0,1)$, $\arg(c+s\tau)$ increases by
$\Arg\tau$ and $\arg(c-s/\tau)$ by $\operatorname{sgn}(\operatorname{Im}\tau)\pi-\Arg\tau$.
Since $\det C=\overline{\det U_X}\det U_Y$ and $D\equiv0\pmod 4$, $\arg\det C$ starts at
$a_0=\sum_k\Arg\tau_k$ and ends at $-a_0+\pi\sum_k\operatorname{sgn}\operatorname{Im}\tau_k
=-a_0+2\pi m$, giving the formula for $r(\pi)$. By~\eqref{eq:inv},
$C(\pi)/r(\pi)=e^{-2\pi i m/D}\,(C(0)/r(0))^{-1}$, so $R(\pi)=(-1)^mR(0)^{\mathsf T}$ and
$P(\pi)=(-1)^{m}\Pf(R(0)^{\mathsf T}-R(0))=(-1)^{m+1}P(0)$.
\end{proof}

\begin{theorem}\label{thm:main}
Let $U\in U(2^n)$, $n\ge3$, have invertible top-row blocks $X,Y$. There is a single-qubit gate $u=R_z(\alpha)R_y(\beta^\ast)$ on
the top qubit such that the block-ZXZ factor $C$ of $Uu$ satisfies~\eqref{eq:split}.
\end{theorem}
\begin{proof}
Rotating by $\alpha$ moves each
$\tau_k$ once around the circle, and the count of $\tau_k$ in the upper half-plane changes by
$\pm1$ at each crossing of the real axis. So there is an open interval of $\alpha$ where no $\tau_k$
is real and the count is even. For such $\alpha$, $m$ is even and Lemma~\ref{lem:parity} gives
$P(\pi)=-P(0)$. The function $\beta\mapsto P(C(\beta),r(\beta))$ is continuous, so it has a zero
$\beta^\ast\in[0,\pi]$ by the intermediate value theorem. At $\beta^\ast$, \eqref{eq:split} holds.
\end{proof}

\begin{corollary}\label{cor:count}
Every $U\in U(2^n)$ is exactly synthesized with
\[
c(n)=4c(n-1)+3\cdot2^{n-1}-6,\quad c(2)=3,\qquad
c(n)=\tfrac{21}{48}4^n-\tfrac32 2^n+2
\]
CNOT gates, which is $c_{\rm ZXZ}(n)-\tfrac{4^{n-2}-1}{3}$.
\end{corollary}
\begin{proof}
At each node apply Theorem~\ref{thm:main} and Lemma~\ref{lem:fold}, and fold two CNOTs of the
input-side multiplexor. All other steps, including the leaf-to-leaf diagonal migration, are as
in block-ZXZ. The gate $u^\dagger$ is a free single-qubit gate. This saves one CNOT per node, and
there are $(4^{n-2}-1)/3$ nodes. For non-generic $U$, note that the circuits produced all share
one CNOT skeleton. The set of unitaries this skeleton realizes (with arbitrary single-qubit gates)
is the continuous image of a compact set, hence closed. It contains the dense set of generic $U$,
so it is all of $U(2^n)$.
\end{proof}

\paragraph{Algorithm (one node).}
(1)~Pick $\alpha$ with an even number of eigenvalues of $e^{i\alpha}X^{-1}Y$ in the upper
half-plane. (2)~Scan $\beta\in[0,\pi]$, tracking the eigenphases of $C(\beta)$ and the root
$r(\beta)$ continuously, locate a sign change of $P$ and refine it by bisection. (3)~Factor
$Uu(\beta^\ast)$ by~\eqref{eq:zxz}, demultiplex $I\oplus C$ with the order and branches of
Lemma~\ref{lem:fold}, fold $\CZ(q_0,t)\CZ(q_{n-2},t)$ into the middle factor, and continue as in
block-ZXZ. Emit $u(\beta^\ast)^\dagger$ on $t$ before the node.

\section{Results}

\begin{table}[h]
\centering
\begin{tabular}{@{}lrrrrrrl@{}}
\toprule
Number of qubits & 1 & 2 & 3 & 4 & 5 & 6 & $n$\\
\midrule
QSD (base)~\cite{smb} & 0 & 6 & 36 & 168 & 720 & 2976 & $(3/4)\cdot4^n-(3/2)\cdot2^n$\\
QSD (optimized)~\cite{smb} & 0 & 3 & 20 & 100 & 444 & 1868 & $(23/48)\cdot4^n-(3/2)\cdot2^n+(4/3)$\\
Block-ZXZ~\cite{zxz} & 0 & 3 & 19 & 95 & 423 & 1783 & $(22/48)\cdot4^n-(3/2)\cdot2^n+(5/3)$\\
\textbf{This work} & \textbf{0} & \textbf{3} & \textbf{18} & \textbf{90} & \textbf{402} & \textbf{1698} & $\boldsymbol{(21/48)\cdot4^n-(3/2)\cdot2^n+2}$\\
\midrule
Theoretical lower bound~\cite{smb} & 0 & 3 & 14 & 61 & 252 & 1020 & $(1/4)\cdot(4^n-3n-1)$\\
\bottomrule
\end{tabular}
\caption{CNOT counts for exact synthesis of a generic $n$-qubit unitary.}
\label{tab:counts}
\end{table}

A reference implementation (numpy; two-qubit leaves by the standard 3-CNOT and 2-CNOT-up-to-diagonal
constructions) reproduces the counts of Table~\ref{tab:counts}. It gives $18$, $90$ and $402$ CNOTs
for Haar-random unitaries at $n=3,4,5$, with reconstruction errors $2\cdot10^{-15}$,
$6\cdot10^{-15}$ and $2\cdot10^{-14}$. The sign relation $P(\pi)=-P(0)$ held on every sampled
target after the choice of $\alpha$, and the root search succeeded on $300/300$ random three-qubit
targets.

\section{Discussion}

\paragraph{Is this an improvement over block-ZXZ?}
Yes, for exact synthesis with guaranteed counts. The leading coefficient drops from
$\tfrac{22}{48}$ to $\tfrac{21}{48}$, about $4.5\%$ fewer CNOTs asymptotically. The count is lower
at every $n\ge3$ (Table~\ref{tab:counts}), and the three-qubit count of $18$ is below the $19$ of
the best previous analytic method. The mechanism is the one block-ZXZ already uses, a CNOT
folded into the middle factor as a CZ. The new ingredient is spending the free top-qubit
single-qubit gates to make a second fold possible. The improvement is modest in absolute terms:
the gap to the lower bound ($\tfrac{12}{48}$) remains, and numerical compilers reach fewer CNOTs
for small $n$~\cite{brickwall,squander}, though without closed forms or guarantees.

\paragraph{Cost and limitations.}
The construction is not closed form: each node needs a one-dimensional root search. As implemented,
evaluating the sign of $P$ ranges over the $\binom{D}{D/2}/2$ balanced splits. That is cheap for
$n\le5$ but grows quickly, and an efficient sign evaluation (or a direct way to find a
split crossing zero) is left open. Existence and the CNOT count do not depend on it.

\paragraph{Outlook.}
The output-side multiplexor has the same fold condition for $A_1A_2^\dagger$, and an $R_y(\pi)$ on the
top qubit at the \emph{output} inverts that spectrum, just as~\eqref{eq:inv} does for $C$.
Numerically, top-qubit gates meeting both conditions exist for every sampled target and give
exact $17$-CNOT three-qubit circuits. A proof would lower the leading coefficient to
$\tfrac{20}{48}$, i.e.\ $\tfrac{20}{48}4^n-\tfrac32 2^n+\tfrac73$ CNOTs.

\begin{thebibliography}{9}
\bibitem{smb} V.~V.~Shende, S.~S.~Bullock, I.~L.~Markov, \emph{Synthesis of quantum logic
circuits}, IEEE TCAD 25, 1000 (2006), arXiv:quant-ph/0406176.
\bibitem{zxz} A.~M.~Krol, Z.~Al-Ars, \emph{Beyond quantum Shannon decomposition: circuit
construction for $n$-qubit gates based on block-ZXZ decomposition}, Phys.\ Rev.\ Applied 22,
034019 (2024), arXiv:2403.13692.
\bibitem{mottonen} M.~M\"ott\"onen, J.~J.~Vartiainen, V.~Bergholm, M.~M.~Salomaa,
\emph{Quantum circuits for general multiqubit gates}, Phys.\ Rev.\ Lett.\ 93, 130502 (2004),
arXiv:quant-ph/0404089.
\bibitem{brickwall} D.~Wierichs, J.~S.~Kottmann, N.~Killoran, \emph{Unitary synthesis with optimal
brick wall circuits}, arXiv:2511.16736.
\bibitem{squander} P.~Rakyta, Z.~Zimbor\'as, \emph{Approaching the theoretical limit in quantum
gate decomposition}, Quantum 6, 710 (2022).
\end{thebibliography}

\end{document}
